\section{Understanding the Modulo Trick}

\begin{frame}{How the Modulo Trick Works - Flowchart}
    \begin{figure}[h!]
        \centering
        \scalebox{0.5}{\begin{tikzpicture}[
            node distance=1cm,
            every node/.style={font=\small},
            box/.style={
                rectangle,
                draw,
                rounded corners,
                align=center,
                minimum width=4.5cm,
                minimum height=0.9cm,
                fill=blue!10
            },
            decision/.style={
                diamond,
                draw,
                align=center,
                aspect=2,
                inner sep=2pt,
                fill=green!10
            },
            arrow/.style={->, thick}
        ]
        
        % Nodes
        \node[box] (eq) {$1+5^a = 7^b + 7^c$};
        
        \node[box, below=of eq] (choosem) {Choose modulus $m$};
        
        \node[box, below=of choosem] (mod) {Find condition on exponents from $1+5^a \equiv 7^b+7^c \pmod m$};
        
        \node[decision, below=1cm of mod] (bound) {Found upper bound\\ of $a$ or $b,c$?};
        
        \node[box, below=of bound] (prop) {Check all finitely many cases};
        
        \node[box, below=of prop] (complete) {Complete solution found};
        
        % Arrows
        \draw[arrow] (eq) -- (choosem);
        \draw[arrow] (choosem) -- (mod);
        \draw[arrow] (mod) -- (bound);
        
        \draw[arrow] (bound) -- node[right] {Yes} (prop);
        \draw[arrow] (prop) -- (complete);
        
        % False loop
        \draw[arrow]
            (bound.east)
            -- node[pos=0.05, right, yshift=7pt] {No}
            ++(2,0)
            |- (choosem.east);
        
        \end{tikzpicture}}
        \caption{The modulo trick method applied to $1+5^a = 7^b+7^c$}
    \end{figure}
\end{frame}

\begin{frame}{Brenner and Foster's Success with Modulo Trick}
    Using the modulo trick, Brenner and Foster successfully solved many equations of the form $1+p^a = q^b + r^c$:
    
    \vspace{0.3cm}
    Examples they solved:
    \begin{itemize}
        \item $1+5^a = 7^b + 7^c$ (as we just saw)
        \item $1+2^a = 3^b + 5^c$
        \item $1+13^a = 17^b + 17^c$
        \item Many others with different base combinations
    \end{itemize}
    
    \vspace{0.3cm}
    The method worked well for these cases - they could always find some modulus that forced an upper bound on at least one variable.
\end{frame}

% ================================================================================
% Section 4: The Problem - Equations That Resist the Modulo Trick
% ================================================================================
